% !TeX root = ../main.tex

\chapter{Related Work}\label{chapter:related-work}

This chapter positions the thesis relative to existing work on travel recommendation systems and recent agentic approaches. The reviewed systems are grouped into interface-oriented travel recommenders and approaches that explicitly incorporate language-model-based reasoning or orchestration.

\section{Travel Recommendation Systems with Classical Interfaces}\label{section:related-work:classical-interfaces}
Existing travel recommendation systems with classical interfaces provide an important baseline for this thesis because they inherit many of their interaction principles from recommender systems more broadly. In the general recommender-systems literature, recommendation has long been operationalized through structured user input, explicit preference models, ranked result lists, filtering controls, and inspectable item representations rather than through open-ended conversation~\parencite{adomaviciusTuzhilinNextGeneration}. This emphasis on usable interaction is also reflected in work that frames recommender systems not only as algorithmic ranking tools but as systems whose value depends on the user experience they provide~\parencite{konstanRiedlUserExperience}. Travel recommendation adopted these same design principles and adapted them to tourism-specific decision tasks. Until the recent spread of large-language-model-based conversational systems, the field was therefore naturally dominated by classical interfaces such as search forms, faceted filters, maps, ranked lists, and stepwise preference elicitation dialogs. A broad survey of tourism recommender systems likewise describes the area primarily through these graphical and form-based interaction patterns and the recommendation pipelines behind them~\parencite{borrasIntelligentTourismSurvey}. This historical dominance followed from the structure of pre-LLM recommender systems: users had to express their needs through explicit attributes or selectable constraints, the system transformed these inputs into recommendation logic, and the resulting candidates were returned in inspectable views that supported comparison and refinement~\parencite{luRecommenderSystems}.

At the algorithmic level, both general recommender systems and travel-specific systems have typically relied on collaborative filtering, content-based recommendation, knowledge-based reasoning, or hybrid combinations of these paradigms~\parencite{burkeHybridSurvey}. Collaborative filtering in particular became one of the foundational paradigms of recommendation, but its effectiveness depends on the availability of suitable interaction data~\parencite{suKhoshgoftaarCollaborativeFiltering}. In the travel domain, knowledge-based and hybrid approaches have been especially relevant because destination choice depends on sparse historical feedback, hard constraints such as budget or season, and subjective preferences that are difficult to infer reliably from behavior alone~\parencite{chaudhariTravelSurvey}. Classical interfaces complement these approaches by making constraints explicit and by enabling iterative preference refinement through filtering, critiquing, and direct inspection of alternatives~\parencite{chenPuCritiquing}. Their main strength is transparency and user control. Their main limitation is that users must adapt nuanced travel goals to the structure of the available controls.

The Destination Finder application illustrates a lightweight form of this interaction style. It centers destination exploration around direct user input and interactive browsing rather than open-ended conversational interpretation~\parencite{destinationFinder}. This type of interface keeps the search space visible and easy to manipulate, but it does not inherently resolve the challenge of interpreting complex, open-ended travel requests.

DestiRec investigates user control in complex recommender systems~\parencite{saricaDestirec}. From the perspective of the present thesis, this line of work is relevant because it frames recommendation not only as an accuracy problem but also as an interaction problem. In travel recommendation, this distinction matters: even technically appropriate suggestions may be less useful if users cannot understand, influence, or iteratively refine the recommendation process. Sarica's thesis is also related to the present work through its use of map-based visualization and a destination dataset in a travel recommendation setting.

\section{Agentic and Language-Model-Based Travel Recommendation Approaches}\label{section:related-work:agentic-approaches}
Recent work no longer treats language models only as generators of final recommendation text. Instead, they are increasingly used as orchestration components that interpret requests, select tools, trigger retrieval, and coordinate downstream recommendation steps. Peng et al. survey this development at the recommender-systems level and describe LLM-powered agents as a shift from static pipelines toward systems that can plan and adapt across multiple steps~\parencite{pengLLMAgentSurvey}. This perspective is directly relevant to the present thesis because the proposed prototype also separates interpretation, retrieval, and response generation rather than relying on one monolithic prompt.

In travel recommendation, this shift is visible in work that combines conversational preference elicitation with data access and explanation. Chen presents an agentic trip-planning application in which planning is structured as a workflow rather than a one-shot answer~\parencite{chenAgenticTripPlanning}. Tandon and Banerjee study intent classification and hybrid retrieval in a conversational tourism recommender, thereby addressing the central problem of mapping natural-language requests to retrieval operations over recommendation data~\parencite{tandonBanerjeeIntentRetrieval}. TRACE extends this line of work toward sustainable tourism by combining conversation with agentic counterfactual explanations, which shows that agentic designs can also be used to justify and gently steer recommendation choices rather than only generate them~\parencite{banerjeeTrace}.

At the same time, these developments are not specific to tourism. Gao et al. describe conversational recommender systems more generally as an interaction setting in which preference elicitation, recommendation, and dialogue management must be coordinated jointly~\parencite{gaoCRSSurvey}. Earlier work by Sun and Zhang already formalized conversational recommendation as a multi-turn decision process, which provides an important conceptual bridge between classical interactive recommenders and newer agentic designs~\parencite{sunZhangCRS}. More recent LLM-based work in e-commerce argues explicitly that conversational recommendation and language models are mutually reinforcing because LLMs can interpret richer user intents while the recommender system provides item grounding and task structure~\parencite{liuEcommerceLLMCRS}.

Another relevant line of work concerns hybrid interaction beyond the travel domain. Feng et al. describe a large-language-model-enhanced conversational recommender system in which the language model is embedded into a broader recommendation stack instead of acting as a standalone chatbot~\parencite{fengLLMEnhancedCRS}. In parallel, research on latent linear critiquing shows that effective recommendation often combines dialogue-like feedback with structured comparison of alternatives, which is closely aligned with hybrid interaction ideas~\parencite{liLatentLinearCritiquing}. More generally, hybrid graphical-conversational interface research argues that conversational and graphical elements should be combined because they support different but complementary user tasks~\parencite{jarvinenPromptingFuture}.

\section{Positioning of This Thesis}\label{section:related-work:positioning}
% TODO to chyba do wyrzucenia, albo cos z tym zrob
The reviewed work suggests two complementary developments. On the one hand, classical travel recommenders emphasize explicit user control and inspectable interaction. On the other hand, recent agentic approaches address intent interpretation, retrieval, and recommendation generation in more flexible conversational settings. This thesis is positioned at the intersection of these directions.

More specifically, the thesis investigates a prototype that combines data-grounded conversational recommendation with a hybrid user interface that exposes recommendation outcomes in a form that users can inspect and refine. Based on the currently available source material, the explicit coupling of multi-turn conversational recommendation, agentic workflow decomposition, and a hybrid interaction layer appears to be a useful focal point for this work. % TODO: refine this positioning after a fuller review of the cited papers.
